\documentclass[11pt]{article}
\usepackage[a4paper,margin=19mm,headheight=14pt]{geometry}
\usepackage[T1]{fontenc}\usepackage{mathpazo}
\usepackage{amsmath,amssymb,enumitem,microtype,fancyhdr}
\usepackage[hidelinks]{hyperref}
\linespread{1.04}\setlength{\parindent}{0pt}\setlength{\parskip}{5pt}
\pagestyle{fancy}\fancyhf{}\fancyhead[L]{\small CCSGD / Reading guide}\fancyfoot[C]{\thepage}\renewcommand{\headrulewidth}{0pt}
\begin{document}
{\LARGE\bfseries CCSGD: a reading guide}\par{\small Aligned English edition | 11 October 2026}\par

\section*{1. What problem does the paper solve?}
The target is the clean objective $F(\theta)=\mathbb E\ell(\theta;Z)$, but only $W=Z+U$ is observed. Even centered input noise generally biases $g(\theta;W)$. Decreasing the learning rate does not remove this population discrepancy. The paper builds an executable stochastic-gradient update for the clean objective.

The main theory assumes Gaussian measurement error with known positive semidefinite covariance, sub-Gaussian latent inputs, and fresh independent samples. Entire extension and input-growth control concern the score, while convexity concerns the objective as a function of the parameter. These are different requirements.


\section*{2. How are the problems classified?}
First separate correction from optimization. Exact Gaussian correction needs an entire input extension with integrability; it needs no convexity. To obtain optimization guarantees, derive a relative variance bound and then select the appropriate objective class.

Convex objectives: Theorem 4.1 bounds the objective gap of a step-size-weighted average. Strongly convex objectives: Theorem 4.2 proves almost-sure and mean-square convergence to $\theta^*$; Theorems 4.3--4.4 give last-iterate and uniformly averaged CLTs under extra local conditions. Nonconvex objectives: Theorem 4.5 controls the weighted mean squared gradient, not distance to a global optimizer.

Then classify the noise. Gaussian noise gives the main theorem package. Uniform and symmetric two-point noise admit explicit conjugate perturbations, discussed in Section 7.1. Their exponential tails restrict admissible score growth. Entire scores with excessive growth and arbitrary non-entire scores are outside the automatic guarantee.


\section*{3. What is the key mechanism?}
Draw $V_t\sim N(0,\Omega)$ independently and set $G_t=\operatorname{Re}g^{\mathbb C}(\theta_t;W_t+iV_t)$, $\theta_{t+1}=\theta_t-\eta_tG_t$. The classical complex corrected-score identity makes each update exactly unbiased. The paper connects this oracle to optimization through explicit moment estimates; no closed-form integration is required.

The key envelope is $\|D_zg^{\mathbb C}(\theta;q)\|\le M(\theta)e^{a\|q\|}$ with $M(\theta)^2\le A+B\|\nabla F(\theta)\|^2$. Conditional orthogonality separates clean sampling noise from correction error, giving $\mathbb E_t\|\xi_t\|^2\le\sigma^2+\rho\|\nabla F(\theta_t)\|^2$. Higher moments are derived from growth and tails, not from a second-moment assumption alone.


\clearpage

\section*{4. Which conclusions are easy to reuse?}
Use $\eta_t=c/(t+\gamma)^\alpha$, $1/2<\alpha\le1$. Convex objective gaps and nonconvex weighted stationarity are $O(T^{\alpha-1})$ for $\alpha<1$, and $O(1/\log T)$ at $\alpha=1$, with the stated step-size cap. Strongly convex MSE is $O(t^{-\alpha})$ for $\alpha<1$; at $\alpha=1$ it is $O(t^{-1})$ when $2\mu c>1$, with logarithmic and slower cases at and below the threshold.

For CLTs, first use convergence, then localize and verify Lindeberg using derived $(2+\delta)$ moments. Covariance continuity and local Hessian regularity are still required. The endpoint $\alpha=1$ needs $c\lambda_{\min}(H)>1/2$ and has a different PR covariance; do not substitute the interior formula $H^{-1}S_*H^{-1}$.

For non-Gaussian conjugates, $a<\lambda$ suffices for the identity and $2a<\lambda$ for the moment reduction, where $\lambda=\pi/b$ for uniform error and $\lambda=\pi/(2b)$ for two-point error. The latter strict inequality also allows a suitable $\delta>0$. Polynomial scores readily satisfy input moment requirements; amplitude and objective assumptions must still be verified.


\section*{5. What do the experiments establish?}
Nonlinear-feature regression is strongly convex in the coefficients and has a separate joint-response moment derivation in Appendix A.9. The parameter MSE slope is $-0.978$, consistent with $O(t^{-1})$. Dashed reference lines are population asymptotic squared biases of the specified updates, not finite-time lower bounds.

The sine-network experiment measures clean-input prediction MSE and is empirical. Fresh measurement noise is drawn at each presentation of fixed latent training pairs; unbiasedness targets the clean empirical objective. No global network-convergence or zero prediction-error limit is claimed. SIMEX has slightly smaller reported MSE at 200,000 updates. WH denotes their deterministic Eq. (5), not all methods in that work.


\section*{6. Recommended reading routes}
For the algorithm: Section 3 and Algorithm 1, then Section 5 examples. For a reusable proof: Lemma 3.4 $\to$ Appendix A.2 $\to$ the relevant theorem in Section 4. For CLTs: Theorem 4.2 $\to$ Appendix A.5 $\to$ A.6 or A.7. For non-Gaussian implementation: Section 7.1 and Algorithm 2. For experimental scope: Section 6, Appendix A.9, and Appendix B.

Agents should start with the five-page AI edition, then resolve stable result IDs in the theorem graph. The original paper is authoritative. The companion conjugate criterion is an external dependency, and the author-provided 6 Astra review reports no major mathematical gaps for main(6).pdf; file-level correspondence with this snapshot is unconfirmed.

\end{document}
